\section{Search in a 2D Matrix}
\label{sec:bs-search-2d-matrix}

\problemstatement{We are given an $n \times m$ matrix with the following
guarantees: every row is sorted in increasing order, and the first element
of every row (except the first) is strictly greater than the last element
of the previous row. Given a target $X$, determine whether $X$ exists in
the matrix, in $O(\log(nm))$ time.}

\begin{patternbox}
The extra guarantee here -- \emph{each row's first element exceeds the
previous row's last element} -- is what elevates this beyond ``each row is
sorted independently'' (Section~\ref{sec:bs-row-max-ones}): it means that
if you read the matrix row by row, left to right, the values you encounter
are \textbf{globally non-decreasing}, exactly as if the matrix were a
single flattened, sorted 1D array. Whenever a 2D matrix problem gives you
this row-chaining guarantee, immediately think ``flatten conceptually and
run ordinary binary search,'' rather than reaching for a more complex 2D
technique.
\end{patternbox}

\subsection*{Understanding the problem carefully}

Because the matrix reads as a single sorted sequence when traversed
row-major (row $0$ left to right, then row $1$ left to right, and so on),
we can treat it as a virtual 1D array of length $nm$, where virtual index
$\textit{idx}$ corresponds to matrix cell
$(\textit{idx} / m,\ \textit{idx} \bmod m)$ (integer division for the row,
remainder for the column). No actual flattening or copying is needed --
we simply convert each binary search midpoint index into a
$(\textit{row}, \textit{col})$ pair on the fly.

\begin{ideabox}[Candidate Space and Predicate]
The candidate space is virtual indices $0, \dots, nm-1$. The predicate is
exactly the same as ordinary binary search
(Section~\ref{sec:bs-find-x}): compare the value at the virtual index to
$X$ and discard the appropriate half.
\end{ideabox}

\subsection*{Why the row-chaining guarantee makes flattening valid}

Ordinary binary search needs the entire candidate sequence (values, in
virtual-index order) to be sorted. The row-chaining guarantee is precisely
what makes this true here: within a row, sortedness gives increasing
values left to right; across the row boundary, the guarantee that the
next row's first element exceeds the previous row's last element ensures
no drop in value occurs when virtual index crosses from the end of one row
to the start of the next. So the entire virtual sequence is non-decreasing,
and ordinary binary search's discard argument
(Section~\ref{sec:bs-find-x}) applies without any modification beyond the
index-mapping arithmetic.

\subsection*{Algorithm}

\begin{enumerate}
  \item Initialize $\textit{lo} \gets 0$, $\textit{hi} \gets nm-1$.
  \item While $\textit{lo} \le \textit{hi}$:
  \begin{enumerate}
    \item $\textit{mid} \gets \textit{lo}+(\textit{hi}-\textit{lo})/2$.
    \item $\textit{row} \gets \textit{mid}/m$, $\textit{col} \gets
          \textit{mid} \bmod m$.
    \item If $\textit{matrix}[\textit{row}][\textit{col}] = X$: return
          \texttt{true}.
    \item If $\textit{matrix}[\textit{row}][\textit{col}] < X$:
          $\textit{lo} \gets \textit{mid}+1$.
    \item Else: $\textit{hi} \gets \textit{mid}-1$.
  \end{enumerate}
  \item If the loop ends without returning, return \texttt{false}.
\end{enumerate}

\begin{algorithm}[H]
\caption{Search in a 2D Matrix}
\begin{algorithmic}[1]
\Procedure{SearchMatrix}{\textit{matrix}, $X$}
  \State $\textit{lo} \gets 0$, $\textit{hi} \gets nm-1$
  \While{$\textit{lo} \le \textit{hi}$}
    \State $\textit{mid} \gets \textit{lo}+(\textit{hi}-\textit{lo})/2$
    \State $\textit{row} \gets \textit{mid} / m$, $\textit{col} \gets \textit{mid} \bmod m$
    \If{$\textit{matrix}[\textit{row}][\textit{col}] = X$} \Return \textbf{true} \EndIf
    \If{$\textit{matrix}[\textit{row}][\textit{col}] < X$}
      \State $\textit{lo} \gets \textit{mid} + 1$
    \Else
      \State $\textit{hi} \gets \textit{mid} - 1$
    \EndIf
  \EndWhile
  \State \Return \textbf{false}
\EndProcedure
\end{algorithmic}
\end{algorithm}

\subsection*{Dry run}

\dryrun{Take $\textit{matrix} = \begin{bmatrix} 1&3&5&7 \\ 10&11&16&20 \\
23&30&34&60 \end{bmatrix}$, $X = 3$. Here $n=3$, $m=4$, $nm=12$.}

\begin{itemize}
  \item $\textit{lo}=0,\textit{hi}=11$: $\textit{mid}=5$,
        $\textit{row}=1,\textit{col}=1$,
        $\textit{matrix}[1][1]=11 > 3$. $\textit{hi}=4$.
  \item $\textit{lo}=0,\textit{hi}=4$: $\textit{mid}=2$,
        $\textit{row}=0,\textit{col}=2$,
        $\textit{matrix}[0][2]=5 > 3$. $\textit{hi}=1$.
  \item $\textit{lo}=0,\textit{hi}=1$: $\textit{mid}=0$,
        $\textit{row}=0,\textit{col}=0$,
        $\textit{matrix}[0][0]=1 < 3$. $\textit{lo}=1$.
  \item $\textit{lo}=1,\textit{hi}=1$: $\textit{mid}=1$,
        $\textit{row}=0,\textit{col}=1$,
        $\textit{matrix}[0][1]=3 = 3$. Return \texttt{true}.
\end{itemize}

\subsection*{C++ implementation}

\begin{lstlisting}
#include <bits/stdc++.h>
using namespace std;

bool searchMatrix(vector<vector<int>>& matrix, int X) {
    int n = matrix.size(), m = matrix[0].size();
    int lo = 0, hi = n * m - 1;

    while (lo <= hi) {
        int mid = lo + (hi - lo) / 2;
        int row = mid / m, col = mid % m;

        if (matrix[row][col] == X) return true;
        else if (matrix[row][col] < X) lo = mid + 1;
        else hi = mid - 1;
    }
    return false;
}

int main() {
    vector<vector<int>> matrix = {
        {1, 3, 5, 7},
        {10, 11, 16, 20},
        {23, 30, 34, 60}
    };
    cout << (searchMatrix(matrix, 3) ? "true" : "false") << endl;
    return 0;
}
\end{lstlisting}

\begin{complexitybox}
\textbf{Time Complexity.} The virtual candidate range of size $nm$ halves
each iteration, giving
\[
  O(\log(nm)).
\]
\textbf{Space Complexity.} $O(1)$ auxiliary space.
\end{complexitybox}

\begin{takeawaybox}
A 2D structure with a strong enough chaining guarantee between rows can be
searched with completely unmodified 1D binary search, as long as you
translate a single virtual index into a $(\textit{row}, \textit{col})$
pair via division and modulo. Always check first whether a 2D problem's
sortedness guarantee is this strong before reaching for a more specialized
2D technique like the one needed in
Section~\ref{sec:bs-search-2d-matrix-2}.
\end{takeawaybox}
